\section{Experimental Design}
\label{sec:design}

\paragraph{Research questions.}
RQ1: does SAT-SA emit the finding families declared for controlled scenarios?
RQ2: how does it compare with simple baselines, and what does each worker
contribute? RQ4: how does it behave on imperfect evidence? RQ5: does verification
detect tampering with the finalized record? RQ6: how sensitive is the peer worker
to cohort size and dispersion? RQ7: what do durable orchestration and review cost,
and do interrupted runs recover? EXT: does the pipeline run on independent real
data, and is its risk associated with an outcome available there? (RQ3, analyst
utility, requires a human study that was not run.) Table~\ref{tab:experiments}
maps questions to experiments.

\widetable{tables/tab-experiments}

\paragraph{Data.}
\emph{Controlled catalog:} \V{C01-action-total} executable scenarios authored
with declared finding families and recommended supervisory actions
(\V{C01-catalog-size} catalog entries in total). Labels are construction labels:
they encode what each scenario was built to contain, so agreement shows the
mechanism responds to its designed inputs, not that it detects real misconduct.
\emph{Generated populations} (PR01): \V{PR01-replicates} populations of
\V{PR01-entities} entities each, \V{PR01-pathological} of them declared
pathological, seeded independently. \emph{Engineered peer cohorts} (P01):
\V{P01-cells} cells crossing cohort size, dispersion and deviation.
\emph{External:} the UCI incident management process event log~\cite{amaral2018uci498}
(dataset 498), a real IT service-desk log of one organisation, mapped to SAT-SA
categories by an adapter and analysed per assignment group
(\V{X02b-groups} groups with enough incidents). It is IT incident data, not SOC
data. Two security datasets were assessed and not used: GUIDE~\cite{freitas2024guide}
has real triage-graded SOC incidents but no acknowledgement, closure,
investigation-step or escalation records (and its download requires an
authenticated account), and CIC-IDS2017~\cite{sharafaldin2018cicids} contains
network flows rather than SOC workflow records.

\paragraph{Baselines.}
For closure-time detection: fixed threshold, MAD rule, z-score and
interquartile-range rules, and random. For prioritization: random order
(\V{PR01-replicates} populations with repeated random trials), critical-and-high alert volume and a
closure-speed heuristic (fastest median closure time first). Externally: random, volume,
reassignment rate and slowest median resolution---the last is close to the
outcome's own construct and is included as an upper reference, not a competitor.

\paragraph{Metrics.}
Precision, recall and F1 against construction labels; precision, recall and
NDCG~\cite{jarvelin2002ndcg} within the top $k\in\{10,20,30\}$\% of entities
(written P@$k$, R@$k$, NDCG@$k$);
review volume needed to find all pathological entities; conformance with the
validation contract and change of emitted finding families relative to a paired control;
wall-clock seconds, database statements and share of time outside the analytics;
Spearman $\rho$ against the SLA-miss rate.

\paragraph{Statistics policy.}
We report medians; 95th percentiles only when $n\ge20$; seeded percentile
bootstrap intervals~\cite{efron1993bootstrap} only when $n\ge10$ independent
units; paired differences with counts of units where one arm was higher, equal or
lower; and no $p$-values or multiple-comparison claims. The independent unit is
stated for each experiment (Table~\ref{tab:experiments}); repeated trials within
one population or one workflow are not independent. Analyses are descriptive.

\paragraph{Evidence discipline.}
Each experiment writes a write-once bundle (manifest with code commit and a
dirty-tree flag, configuration, raw results, processed metrics, summary) with
SHA-256 hashes. The canonical bundles are frozen in
\texttt{research/evidence/freeze-v2}; superseded bundles (for example X02,
replaced by X02b after an ingestion-bookkeeping correction that left every
analytical output identical) are kept but not cited for results.
Every number in this paper is rendered from these bundles by the paper data layer
(Section~\ref{sec:repro}).
